\begin{statementbox}

	\textbf{\textcolor{CStatement}{条件}}
	\begin{description}[style=nextline, leftmargin=2.8em, labelsep=0.5em, itemsep=0.35em]
		\item[\textbf{\textcolor{CStatement}{条件 1（互不相容）：}}] $A_i\cap A_j=\varnothing\;(i\neq j)$.
		\item[\textbf{\textcolor{CStatement}{条件 2（并集完备）：}}] $\bigcup_{i=1}^n A_i=\Omega$.
		\item[\textbf{\textcolor{CStatement}{条件 3（概率正性）：}}] $P(A_i)>0$ 对每个 $i$.
	\end{description}

	\textbf{\textcolor{CStatement}{结论}}
	\begin{description}[style=nextline, leftmargin=2.8em, labelsep=0.5em, itemsep=0.45em]
		\item[\textbf{\textcolor{CStatement}{结论一（全概率公式）：}}]
			\textbf{\textcolor{CStatement}{直观描述：}}
			将总概率按原因分类加权求和。
			\[
				P(B)=\sum_{i=1}^{n}P(A_i)P(B\mid A_i).
			\]

		\item[\textbf{\textcolor{CStatement}{常见变形（两事件划分）：}}]
			\textbf{\textcolor{CStatement}{直观描述：}}
			样本空间仅分为两个对立部分时的形式。
			\[
				P(B)=P(A)P(B\mid A)+P(\overline{A})P(B\mid \overline{A}).
			\]
	\end{description}

	\textbf{\textcolor{CStatement}{适用范围：}}
	仅当事件组 $\{A_i\}$ 构成 $\Omega$ 的一个划分且每个 $P(A_i)>0$ 时才能使用；若 $A_i$ 出现零概率事件则公式失效。

\end{statementbox}
